\chapter{2023:Moungi Bawendi and Louis Brus and Aleksey Yekimov}

\label{chap:chemistry2023}

The Nobel Prize in Chemistry for the year 2023 was awarded jointly to Moungi Bawendi and Louis Brus and Aleksey Yekimov.

\textbf{Motivation:}

for the discovery and synthesis of quantum dots

\section{Moungi Bawendi}

\subsection*{Early Life and Education}

Moungi Bawendi was born in 1961 in Paris, France.

\subsection*{Career and Major Research}

Bawendi earned a bachelor's degree from Harvard University in 1982 and a doctorate in chemistry from the University of Chicago in 1988. He became a professor at the Massachusetts Institute of Technology, where in 1993 he developed a method for producing near-identical, size-controlled quantum dots in large quantities.

\subsection*{Nobel Prize-winning Work}

The work for which Moungi Bawendi was awarded the Nobel Prize in Chemistry in 2023 was described by the Nobel Committee as: "for the discovery and synthesis of quantum dots".

Bawendi's synthesis of monodisperse quantum dots made the particles available in precise, reproducible sizes, so that their color could be controlled by size alone.

\subsection*{Significance and Impact}

His method turned quantum dots from an object of study into a practical material, enabling their use in displays, lasers, and biomedical imaging.

\subsection*{Later Career and Honors}

Bawendi continues his research on semiconductor nanocrystals at the Massachusetts Institute of Technology.

\subsection*{Legacy}

The synthetic chemistry he developed is the basis of the commercial quantum dot technology used in televisions and screens.

\subsection*{Selected Publications}

"Synthesis and Characterization of Nearly Monodisperse CdE (E = S, Se, Te) Semiconductor Nanocrystallites" (Journal of the American Chemical Society, 1993)

\clearpage

\section{Louis Brus}

\subsection*{Early Life and Education}

Louis Brus was born in 1943 in Cleveland, Ohio, USA.

\subsection*{Career and Major Research}

Brus earned a bachelor's degree from Rice University in 1965 and a doctorate from Columbia University in 1969. He worked at Bell Laboratories and later became a professor at Columbia University. In the early 1980s he discovered that the optical properties of colloidal semiconductor particles change with their size, the quantum size effect now known as quantum confinement.

\subsection*{Nobel Prize-winning Work}

The work for which Louis Brus was awarded the Nobel Prize in Chemistry in 2023 was described by the Nobel Committee as: "for the discovery and synthesis of quantum dots".

Brus showed that in nanometer-sized semiconductor crystallites, the energy levels and color of light absorption depend on the particle size, the defining property of quantum dots.

\subsection*{Significance and Impact}

His discovery established that the same material emits different colors depending on its size, and it founded the chemistry of colloidal quantum dots.

\subsection*{Later Career and Honors}

Brus continues his research on nanocrystals and their spectroscopy at Columbia University.

\subsection*{Legacy}

The size-dependence that Brus explained is the basis of quantum dot technology, from displays to biological imaging.

\subsection*{Selected Publications}

"Electron-Electron and Electron-Hole Interactions in Small Semiconductor Crystallites" (Journal of Chemical Physics, 1984)

\clearpage

\section{Aleksey Yekimov}

\subsection*{Early Life and Education}

Aleksey Yekimov was born in 1945 in Saint Petersburg, Russia.

\subsection*{Career and Major Research}

Yekimov worked as a physicist at the Vavilov State Optical Institute in Leningrad. In 1981 he produced copper chloride nanocrystals in molten glass and showed that their absorption of light depended on their size, demonstrating the quantum size effect in a semiconductor.

\subsection*{Nobel Prize-winning Work}

The work for which Aleksey Yekimov was awarded the Nobel Prize in Chemistry in 2023 was described by the Nobel Committee as: "for the discovery and synthesis of quantum dots".

Yekimov's experiments provided the first clear evidence that the electronic properties of semiconductor nanocrystals change with their dimensions, the phenomenon underlying quantum dots.

\subsection*{Significance and Impact}

His observation in glass-hosted crystals was the starting point for the entire field of size-tunable semiconductor nanoparticles.

\subsection*{Later Career and Honors}

Yekimov worked at the Vavilov State Optical Institute in Leningrad and, since 1999, at Nanocrystals Technology in the United States.

\subsection*{Legacy}

The discovery Yekimov made in glass established quantum dots as a new class of material whose properties are set by size.

\subsection*{Selected Publications}

"Quantum Size Effect in Semiconductor Microcrystals" (Solid State Communications, 1985)

\clearpage

\section*{Chapter Summary}

The 2023 prize recognized the discovery and synthesis of quantum dots: Yekimov first observed size-dependent properties in glass, Brus explained the quantum size effect in colloidal particles, and Bawendi developed the synthesis of uniform nanocrystals. Quantum dots are now used in displays and biomedical imaging.
